% TOP_PROBLEM: {"id": 2808, "title": "Kirby Problem 3.10", "category_id": 7, "category": {"id": 7, "name": "topology", "display_name": "Topology", "description": "Properties preserved under continuous deformations.", "slug": "topology", "order_index": 7, "created_at": "2026-07-31T15:26:25.671Z"}, "status": "open", "problem_number": "KP-3.10", "set_id": 11}
% FIRST_POSTED: 2026-10-01
% ENTRY_KIND: statement_only
% UnsolvedMath ID 2808; source catalogue code KP-3.10
% !TeX program = lualatex
% Definition and sources only; this file is not an LLM solution attempt.
\documentclass[11pt,a4paper]{article}
\usepackage[margin=25mm]{geometry}
\usepackage{fontspec}
\setmainfont{lmroman10-regular.otf}[BoldFont=lmroman10-bold.otf,
 ItalicFont=lmroman10-italic.otf,BoldItalicFont=lmroman10-bolditalic.otf]
\usepackage{amsmath,amssymb,mathtools}
\usepackage{unicode-math}
\setmathfont{latinmodern-math.otf}
\AtBeginDocument{\let\setminus\smallsetminus}
\usepackage{xurl,hyperref}
\hypersetup{colorlinks=true,urlcolor=blue}
\setcounter{secnumdepth}{0}
\setlength{\parindent}{0pt}
\setlength{\parskip}{5pt}
\setlength{\emergencystretch}{3em}
\begin{document}
\section*{Kirby Problem 3.10}\label{2808}
{\small UnsolvedMath ID 2808 · KP-3.10 · Topology · Kirby's Problems in Low-Dimensional Topology}
\subsection{Definitions and mathematical statement}
A closed connected oriented three-manifold $M$ is a rational homology sphere
if $H_*(M;\mathbb Q)\cong H_*(S^3;\mathbb Q)$, and an integral homology sphere
if the analogous statement holds with $\mathbb Z$ coefficients.
For the latter, in particular, $H_1(M;\mathbb Z)=0$.

An orientable hyperbolic three-manifold is arithmetic if it is
$\mathbb H^3/\Gamma$, with $\Gamma$ a torsion-free lattice in
$\operatorname{PSL}_2(\mathbb C)$ commensurable with the norm-one group of
an order in a quaternion algebra over a number field with exactly one complex
place, ramified at every real place. An order is a full integral lattice which
is also a subring. Commensurability of lattices allows conjugation and means
that their intersection has finite index in each. For manifolds this is
equivalent to possessing a common finite-sheeted cover.

(a) Are there infinitely many commensurability classes containing an
arithmetic hyperbolic rational homology three-sphere?

(b) Are there infinitely many arithmetic hyperbolic integral homology
three-spheres, up to isometry?

The quantifiers differ: an infinite sequence of covers in one commensurability
class would not settle (a), whereas (b) asks for infinitely many distinct
manifolds and does not demand pairwise incommensurability.
\subsection{Short English statement}
Find infinitely many arithmetic hyperbolic rational homology sphere classes, and determine whether there are infinitely many arithmetic integral homology spheres.
\subsection{Sources}
\begin{itemize}
\item[S1] ULAM.AI and the UnsolvedMath Contributors, supplied problems.json, record 2808 (KP-3.10), Kirby's Problems in Low-Dimensional Topology. Catalogue identity and source transcription; CC BY 4.0.
\url{https://huggingface.co/datasets/ulamai/UnsolvedMath}.
\item[S2] K3: A New Problem List in Low-Dimensional Topology, Problem 3.10. Expanded formulation of the supplied catalogue statement.
\url{https://math.berkeley.edu/sites/default/files/surv-295-ruberman-watermarked-author-pdf.pdf}.
\end{itemize}
\end{document}
